% =============================================================================
%  Resumen Capítulo 6 — ANÁLISIS DE AGRUPAMIENTOS
%  Libro: Evsukoff, A. (2020). Inteligência Computacional: Fundamentos e
%         Aplicações. Río de Janeiro: E-papers.
% =============================================================================
% =============================================================================
%  preambulo.tex — Preámbulo de los resúmenes en español
%  Evsukoff, A. (2020). Inteligência Computacional. Rio de Janeiro: E-papers.
% =============================================================================
\documentclass[11pt,a4paper]{article}

\usepackage[T1]{fontenc}
\usepackage[utf8]{inputenc}
\usepackage[spanish,es-tabla]{babel}
\usepackage{lmodern}
\usepackage[a4paper,margin=2.4cm]{geometry}
\usepackage{amsmath,amssymb,amsthm,mathtools}
\usepackage{bm}
\usepackage{enumitem}
\usepackage{xcolor}
\usepackage{tcolorbox}
\tcbuselibrary{skins,breakable}
\usepackage{hyperref}
\usepackage{microtype}
\usepackage{titlesec}
\usepackage{fancyhdr}
\usepackage{booktabs}
\usepackage{array}

\hypersetup{
  colorlinks=true,
  linkcolor=black!70,
  urlcolor=blue!60!black,
  citecolor=black!70,
  pdfborder={0 0 0}
}

\titleformat{\section}{\Large\bfseries\sffamily\color{black!85}}{\thesection}{0.7em}{}
\titleformat{\subsection}{\large\bfseries\sffamily\color{black!75}}{\thesubsection}{0.6em}{}

\pagestyle{fancy}
\fancyhf{}
\fancyhead[L]{\small\itshape Resumen — Inteligencia Computacional (Evsukoff, 2020)}
\fancyhead[R]{\small\thepage}
\renewcommand{\headrulewidth}{0.3pt}

\newtcolorbox{conceito}[1][]{
  colback=blue!4, colframe=blue!55!black, boxrule=0.6pt,
  arc=2mm, left=6pt, right=6pt, top=4pt, bottom=4pt,
  title={\bfseries Concepto clave}, fonttitle=\sffamily\bfseries,
  breakable, #1
}
\newtcolorbox{formula}[1][]{
  colback=gray!6, colframe=gray!50, boxrule=0.4pt,
  arc=1.5mm, left=6pt, right=6pt, top=3pt, bottom=3pt,
  breakable, #1
}
\newtcolorbox{aplicacao}[1][]{
  colback=green!4, colframe=green!45!black, boxrule=0.5pt,
  arc=2mm, left=6pt, right=6pt, top=3pt, bottom=3pt,
  title={\bfseries Aplicación}, fonttitle=\sffamily\bfseries,
  breakable, #1
}

\newcommand{\R}{\mathbb{R}}
\newcommand{\E}{\mathbb{E}}
\newcommand{\Var}{\operatorname{Var}}
\newcommand{\Cov}{\operatorname{Cov}}
\newcommand{\diag}{\operatorname{diag}}
\newcommand{\argmin}{\operatorname*{arg\,min}}
\newcommand{\argmax}{\operatorname*{arg\,max}}
\newcommand{\T}{^{\!\top}}
\DeclareMathOperator{\sign}{sign}
\DeclareMathOperator{\tr}{tr}


\title{\sffamily\bfseries Capítulo 6 — Análisis de Agrupamientos\\
       \large Particionales, jerárquicos, espectrales y validación}
\author{Resumen de estudio}
\date{}

\begin{document}
\maketitle
\thispagestyle{fancy}

\begin{conceito}
El \textbf{análisis de agrupamientos} (\emph{clustering}) divide el
conjunto de datos $\mathcal{T}=\{\mathbf{x}(t)\}_{t=1}^{N}$ en
$K$~subconjuntos disjuntos
$\mathbb{C}=\{\mathcal{C}_1,\dots,\mathcal{C}_K\}$, de modo que los
registros del mismo grupo sean \emph{similares entre sí} y
\emph{disímiles} respecto de los demás grupos. Es un problema
\textbf{no supervisado}: no hay variable-objetivo, lo que hace
intrínsecamente difícil su evaluación. Incluso un conjunto pequeño
admite una explosión combinatoria de particiones ($N=100$ registros en
dos grupos de~50 producen $\sim 10^{29}$ particiones), y la ``mejor''
partición depende siempre del contexto de la aplicación.
\end{conceito}

\section{El problema}

Una partición satisface
$\bigcup_{i}\mathcal{C}_i=\mathcal{T}$ y
$\mathcal{C}_i\cap\mathcal{C}_j=\emptyset$ para $i\neq j$. La
\textbf{forma de los grupos} la determina la métrica de distancia
elegida (Figura~6.3): la Euclidiana ($L_2$) genera esferas; la
Mahalanobis, elipsoides; Manhattan ($L_1$) y $L_\infty$, hipercubos.
Los grupos \emph{convexos} (separables por hiperplanos) los obtienen
métodos particionales simples; los grupos \emph{no convexos} requieren
métodos jerárquicos del tipo \emph{single-link} o métodos espectrales.

\textbf{Atención:} los grupos hallados \emph{no} tienen por qué
coincidir con clases pre-existentes. ``Clasificación no supervisada''
es una metáfora útil pero peligrosa.

\section{Métodos de partición: $k$-medias}

Se minimiza la suma de distancias cuadráticas a los centros:
\begin{formula}
\[
J(\boldsymbol\Omega)=\sum_{t=1}^{N}\sum_{i=1}^{K}
u_i(\mathbf{x}(t))\,\|\mathbf{x}(t)-\boldsymbol\omega_i\|^2,
\quad \boldsymbol\omega_i=\frac{1}{\hat N_i}\!\!\sum_{\mathbf{x}(t)\in\mathcal{C}_i}\!\!\mathbf{x}(t),
\]
\end{formula}
con $u_i(\mathbf{x})=1$ si $\mathbf{x}\in\mathcal{C}_i$, 0 en caso
contrario. Esto equivale a minimizar $\tr(\boldsymbol\Sigma_W)$ y, como
$\boldsymbol\Sigma_T=\boldsymbol\Sigma_W+\boldsymbol\Sigma_B$ es
constante, \textbf{maximizar simultáneamente}
$\tr(\boldsymbol\Sigma_B)$ --- es decir, grupos compactos y bien
separados.

\paragraph{Algoritmo (Lloyd/MacQueen).}
\begin{enumerate}[leftmargin=1.2cm]
  \item Inicializar $\boldsymbol\Omega=[\boldsymbol\omega_1,\dots,\boldsymbol\omega_K]$
        (aleatorio, \emph{k-means++}, PCA-Part, CCIA\dots).
  \item Asignar cada $\mathbf{x}(t)$ al centro más próximo
        ($\argmin_j\|\mathbf{x}(t)-\boldsymbol\omega_j\|$).
  \item Recalcular cada centro como media de los registros asignados.
  \item Iterar hasta $\|\boldsymbol\Omega-\boldsymbol\Omega_0\|<\mathrm{tol}$.
\end{enumerate}

Complejidad $O(NKp)$ por iteración; muy eficiente cuando $Kp\ll N$.

\paragraph{Limitaciones.}
Problema NP-difícil, heurística \emph{voraz}: cae en mínimos locales.
Sensible a (i) inicialización, (ii) presencia de \emph{outliers}, (iii)
grupos de tamaños desiguales o no convexos. Cuando el número de grupos
coincide con la estructura real, las inicializaciones convergen a la
misma solución; cuando no, soluciones distintas en cada ejecución
indican que $K$ está mal elegido.

\paragraph{Variantes.} ISODATA (ajusta $K$ automáticamente),
$k$-medoids/PAM/CLARA/CLARANS (centro = registro real, robusto a
\emph{outliers}), $k$-medians, $x$-means (escoge $K$ por AIC/BIC),
$k$-medias esférico (similitud del coseno, alta dimensión),
\emph{fuzzy} $c$-medias (Capítulo~8).

\section{Métodos jerárquicos}

Construyen una cadena de particiones
$\mathbb{C}_0\subseteq\mathbb{C}_1\subseteq\cdots\subseteq\mathbb{C}_{N-1}$
visualizada como árbol binario (\textbf{dendrograma}). Los
\textbf{aglomerativos} (más comunes) inician con $N$ \emph{singletons}
y fusionan el par más próximo; los \textbf{divisivos} hacen el camino
inverso.

La función de \textbf{ligadura} entre grupos define el resultado:
\begin{itemize}[leftmargin=1.2cm]
  \item \emph{Single-link} (vecino más próximo):
        $L_{\min}(\mathcal{C}_i,\mathcal{C}_j)=\min\,\mathrm{dist}(\boldsymbol\nu,\boldsymbol\upsilon)$
        --- captura grupos no convexos, sensible a ``cadenas''.
  \item \emph{Complete-link} (vecino más lejano):
        $L_{\max}=\max\,\mathrm{dist}$ --- grupos compactos.
  \item \emph{UPGMA} (promedio):
        $L_{\mathrm{med}}=\frac{1}{|\mathcal{C}_i||\mathcal{C}_j|}\sum\sum\mathrm{dist}$.
  \item \emph{UPGMC} (centroides):
        $L_{\mathrm{cen}}=\mathrm{dist}(\bar{\boldsymbol\nu},\bar{\boldsymbol\upsilon})$.
  \item \emph{Ward}: minimiza la varianza intragrupos en cada fusión.
\end{itemize}
El número adecuado de grupos se estima por la mayor distancia entre
dos niveles sucesivos del dendrograma. Costo de memoria $O(N^2)$ por
la matriz de distancias --- inviable para $N$ muy grande.

\section{Particionamiento de grafos y métodos espectrales}

Se representa $\mathcal{T}$ como grafo $G(\mathcal{V},\mathcal{E})$ con
$\mathcal{V}\equiv\mathcal{T}$ y matriz de adyacencias $\mathbf{A}$
ponderada por la similitud. Función usual: gaussiana
$\mathrm{sim}_G(\mathbf{x}(i),\mathbf{x}(j))=\exp(-\|\mathbf{x}(i)-\mathbf{x}(j)\|^2/2\sigma^2)$
o coseno (texto).

\paragraph{Corte mínimo.} Para partición
$\{\mathcal{C}_1,\mathcal{C}_2\}$,
$\mathrm{cut}(\mathcal{C}_1,\mathcal{C}_2)=\frac12\sum_{i,j}a_{ij}u_1(i)u_2(j)$.
Solución trivial: aislar el vértice de menor grado (\emph{outlier}).
El \textbf{corte normalizado de Shi--Malik} corrige:
\[
\mathrm{Ncut}=\frac{\mathrm{cut}(\mathcal{C}_1,\mathcal{C}_2)}{\mathrm{cut}(\mathcal{C}_1,\mathcal{V})}
+\frac{\mathrm{cut}(\mathcal{C}_1,\mathcal{C}_2)}{\mathrm{cut}(\mathcal{C}_2,\mathcal{V})}.
\]

\paragraph{Laplaciano.} $\mathbf{L}=\mathbf{D}-\mathbf{A}$ con
$\mathbf{D}=\diag(d_1,\dots,d_N)$, $d_i=\sum_j a_{ij}$. Es semi-definido
positivo; tiene autovalor $\lambda_1=0$ con multiplicidad igual al
número de componentes conexas. El segundo autovector (\emph{Fiedler})
revela el mejor corte.

\paragraph{Algoritmo de agrupamiento espectral.}
\begin{enumerate}[leftmargin=1.2cm]
  \item Construir $\mathbf{A}$ con la función de similitud.
  \item Calcular un laplaciano: $\mathbf{L}=\mathbf{D}-\mathbf{A}$
        (corte mínimo), generalizado
        $\mathbf{L}\mathbf{W}=\mathbf{D}\mathbf{W}\boldsymbol\Lambda$
        (Shi--Malik), o
        $\hat{\mathbf{A}}=\mathbf{D}^{-1/2}\mathbf{A}\mathbf{D}^{-1/2}$
        (Ng--Jordan--Weiss).
  \item Descomposición espectral; tomar los $K$ autovectores asociados
        a los $K$ menores autovalores (o $K$ mayores en NJW).
  \item Aplicar $k$-medias a la matriz $\mathbf{Z}$ formada por esos
        autovectores (en NJW, normalizar cada fila).
\end{enumerate}
La transformación no lineal (similitud gaussiana) permite separar
grupos no convexos en el espacio original. La inicialización del
$k$-medias se vuelve determinista (matriz identidad). Costo dominado
por $O(N^3)$ de la descomposición.

\paragraph{Modularidad (Newman--Girvan).} Mide cuánto se aleja la
partición de la esperada en una red aleatoria:
\[
Q=\frac{1}{2E}\sum_{ij}\Bigl(a_{ij}-\frac{d_id_j}{2E}\Bigr)\mathbf{u}_i\mathbf{u}_j\T,
\quad \mathbf{B}=\mathbf{A}-\mathbf{d}\T\mathbf{d}/(2E).
\]
Maximizar $Q$ es el problema central de la \emph{detección de
comunidades} en redes complejas; se resuelve por descomposición
espectral de $\mathbf{B}$ o por enfoques jerárquicos divisivos.

\section{Validación de grupos}

\subsection{Externas} (requieren rótulos de clases)
\begin{itemize}[leftmargin=1.2cm]
  \item \textbf{Pureza:}
        $\mathrm{pur}(\mathcal{C}_j)=\max_i n_{ij}/\hat N_j$
        (homogeneidad);
        $\mathrm{PUR}=\sum_j(\hat N_j/N)\mathrm{pur}(\mathcal{C}_j)$.
  \item \textbf{Entropía} (Shannon):
        $\mathrm{ent}(\mathcal{C}_j)=-\frac{1}{\log_2 m}\sum_i (n_{ij}/\hat N_j)\log_2(n_{ij}/\hat N_j)$
        --- menor es mejor.
\end{itemize}

\subsection{Internas} (sólo usan datos y partición)
\begin{itemize}[leftmargin=1.2cm]
  \item \textbf{PI (Bensaid):}
        $\tr(\boldsymbol\Sigma_W)/\sum_{ij}N_i\|\boldsymbol\omega_i-\boldsymbol\omega_j\|^2$.
  \item \textbf{VAT:} \emph{data image} reordenado --- diagnóstico
        visual cualitativo.
\end{itemize}

\subsection{Relativas} (elección de $K$)
\begin{itemize}[leftmargin=1.2cm]
  \item \textbf{Calinski--Harabasz:}
        $\mathrm{CH}=\frac{N-K}{K-1}\frac{\tr(\boldsymbol\Sigma_B)}{\tr(\boldsymbol\Sigma_W)}$ (max).
  \item \textbf{PBM:}
        $(\frac{1}{K}\frac{E_1}{E_K}\frac{1}{D_{\max}})^2$ (max).
  \item \textbf{Silueta:} $sw_i(t)=(b_i(t)-a_i(t))/\max(a_i,b_i)$, con
        $a_i$ distancia media intragrupo y $b_i$ menor media
        entregrupos. $\mathrm{SW}\in(-1,1)$, próximo a~1 es ideal.
  \item \textbf{Xie--Beni:}
        $\tr(\boldsymbol\Sigma_W)/(N\cdot D_{\min})$ (min).
\end{itemize}

Vendramin, Campello \& Hruschka (estudio con 1\,080~datasets) concluyen
que \textbf{CH, PBM y SW} son los índices más efectivos para descubrir
$K$. Se recomienda ejecutar el $k$-medias múltiples veces para cada
$K$ y analizar la dispersión de los índices vía \emph{box-plot}.

\section{Aplicaciones}

\begin{aplicacao}
\textbf{Segmentación de imágenes SAR de la Amazonía} (Beisl,
Coppe/UFRJ). Imágenes del satélite JERS-1 en las épocas de creciente
(oct./1995) y de seca (mayo/1996) sobre la cuenca del Solimões.
Pre-procesamiento por función \emph{semivariograma} caracteriza la
textura local. El agrupamiento por ISODATA (variante de $k$-medias)
seguido de clasificación por especialista produce cuatro clases
temáticas: agua, bosque seco, bosque inundado y vegetación inundada.
El resultado alimenta mapas de sensibilidad ambiental al derrame de
petróleo (poliducto Coari--Manaus, Petrobras).
\end{aplicacao}

\begin{aplicacao}
\textbf{Agrupamiento de documentos} (Santos, Coppe/UFRJ).
585~resúmenes de tesis del PEC/Coppe (2004--2008); el
pre-procesamiento (\emph{stopwords}, \emph{stemming}) reduce a
265~términos; ponderación \emph{tf-idf}. Se construye un grafo con
similitud del coseno y umbral~$\varepsilon$; se aplica el algoritmo de
\textbf{maximización de la modularidad} de Newman--Girvan con etapa de
refinamiento. En $\varepsilon=0{,}20$, $Q=0{,}5334$ y 7~comunidades
coherentes con las áreas de investigación del Programa: Mecánica
Computacional, Estructuras, Geotecnia, Petróleo (Geoquímica y
Geofísica), Sistemas Computacionales y Recursos Hídricos.
\end{aplicacao}

\section{Comentarios finales}

\begin{enumerate}[leftmargin=1.2cm]
  \item El \emph{clustering} es un \textbf{problema mal planteado}:
        ninguna medida es absoluta, el criterio final es siempre
        interpretativo.
  \item \textbf{$k$-medias y jerárquicos aglomerativos} son los
        enfoques clásicos, en cualquier software. Los métodos
        \textbf{espectrales} brillan con grupos no convexos y en alta
        dimensión (textos, imágenes).
  \item \textbf{No confundir grupos con clases}: la coincidencia es
        circunstancial. Los métodos \emph{semi}-supervisados usan
        rótulos parciales para guiar el agrupamiento.
  \item En minería de textos y bioinformática, la \textbf{alta
        dimensionalidad} obliga a usar similitudes específicas
        (coseno) y técnicas de reducción (LSI, ACP).
  \item Próximos capítulos: \emph{fuzzy} $c$-medias (Capítulo~8),
        funciones de base radial (Capítulo~7), redes neuronales
        (Capítulo~9) amplían aún más el arsenal.
\end{enumerate}

\end{document}
